
\documentclass[11pt,a4paper]{article}

\usepackage[utf8]{inputenc}
\usepackage[T1]{fontenc}
\usepackage[french]{babel}
\usepackage{lmodern}
\usepackage{microtype}
\usepackage[a4paper,margin=1.65cm,headheight=15pt]{geometry}
\usepackage{amsmath,amssymb,mathtools}
\usepackage{array,tabularx,multicol}
\usepackage{enumitem}
\usepackage[most]{tcolorbox}
\usepackage{xcolor}
\usepackage{tikz}
\usetikzlibrary{arrows.meta,calc}
\usepackage{pdflscape}
\usepackage{fancyhdr}
\usepackage{fancyvrb}
\usepackage{listings}
\usepackage{hyperref}

\definecolor{fondpage}{RGB}{247,248,250}
\definecolor{encre}{RGB}{31,35,40}
\definecolor{bleu}{RGB}{40,91,135}
\definecolor{vert}{RGB}{55,125,92}
\definecolor{orange}{RGB}{178,105,42}
\definecolor{violet}{RGB}{101,77,145}
\definecolor{rouge}{RGB}{170,72,72}
\definecolor{gris}{RGB}{86,91,99}
\definecolor{bleufonce}{RGB}{28,55,120}
\definecolor{bleuclair}{RGB}{238,243,255}
\definecolor{grisclair}{RGB}{245,245,245}
\definecolor{tableline}{RGB}{45,45,45}
\definecolor{softgray}{RGB}{246,247,249}
\definecolor{deepblue}{RGB}{25,62,115}
\definecolor{softblue}{RGB}{235,241,250}
\definecolor{accent}{RGB}{20,82,130}

\hypersetup{colorlinks=true,linkcolor=bleufonce,urlcolor=bleufonce,pdftitle={Parcours de révision -- Suites -- Terminale spécialité},pdfauthor={}}

\setlength{\parindent}{0pt}
\setlength{\parskip}{0.25em}
\renewcommand{\arraystretch}{1.12}
\setlength{\tabcolsep}{3pt}
\setlist[itemize]{leftmargin=1.25em,itemsep=0.15em,topsep=0.2em}
\setlist[enumerate,1]{label=\textbf{\arabic*.}, leftmargin=1.05cm, itemsep=0.3em, topsep=0.2em}
\setlist[enumerate,2]{label=\textbf{\alph*.}, leftmargin=0.85cm, itemsep=0.2em, topsep=0.15em}

\pagestyle{fancy}
\fancyhf{}
\lhead{}
\rhead{}
\cfoot{\thepage}

\newcommand{\R}{\mathbb{R}}
\newcommand{\N}{\mathbb{N}}
\newcommand{\Un}{\ensuremath{(u_n)}}
\newcommand{\e}{\mathrm e}
\newcommand{\vect}[1]{\overrightarrow{#1}}
\newcommand{\origine}[1]{ {\scriptsize\textcolor{bleufonce}{(site #1)}}}

\newcounter{question}
\newcounter{reponse}
\newenvironment{blocquestions}[1]{%
  \begin{tcolorbox}[enhanced,breakable,colback=bleuclair,colframe=bleufonce,boxrule=0.6pt,arc=2mm,left=2mm,right=2mm,top=2mm,bottom=1.5mm,title={\bfseries #1},coltitle=white,colbacktitle=bleufonce,before skip=3mm,after skip=3mm, fontupper=\large]
  \large
  \begin{enumerate}[leftmargin=*,label=\textbf{\arabic*.},itemsep=0.85em,topsep=0.5em]
  \setcounter{enumi}{\value{question}}
}{%
  \setcounter{question}{\value{enumi}}
  \end{enumerate}
  \end{tcolorbox}
}
\newenvironment{blocreponses}[1]{%
  \begin{tcolorbox}[enhanced,breakable,colback=bleuclair,colframe=bleufonce,boxrule=0.6pt,arc=2mm,left=2mm,right=2mm,top=2mm,bottom=1.5mm,title={\bfseries #1},coltitle=white,colbacktitle=bleufonce,before skip=3mm,after skip=3mm, fontupper=\large]
  \large
  \begin{enumerate}[leftmargin=*,label=\textbf{\arabic*.},itemsep=0.45em,topsep=0.5em]
  \setcounter{enumi}{\value{reponse}}
}{%
  \setcounter{reponse}{\value{enumi}}
  \end{enumerate}
  \end{tcolorbox}
}

\newtcolorbox{correctionbox}[1]{enhanced,breakable,colback=grisclair,colframe=black!55,boxrule=0.55pt,arc=2mm,left=2mm,right=2mm,top=2mm,bottom=1.5mm,title=\textbf{#1},coltitle=black,colbacktitle=black!12,before skip=2mm,after skip=2mm}

\newcommand{\titreSujet}[2]{%
  \subsection*{#1}
  \addcontentsline{toc}{subsection}{#1}
  \textit{Référence : #2}\par\medskip\hrule\medskip
}

\newcommand{\sep}{\vspace{0.55mm}}
\tcbset{enhanced,arc=2mm,outer arc=2mm,boxrule=0.42pt,left=1.15mm,right=1.15mm,top=0.75mm,bottom=0.75mm,boxsep=0.35mm,before skip=0.8mm,after skip=0.8mm,fonttitle=\bfseries\footnotesize,coltitle=encre,before upper=\raggedright\fontsize{7.65}{8.15}\selectfont}
\newtcolorbox{fichebox}[3][]{colback=white,colframe=#2!72!black,colbacktitle=#2!18,coltitle=#2!50!black,borderline west={0.9mm}{0pt}{#2!78!black},title={#3},drop fuzzy shadow=black!5,#1}
\newcommand{\tagcol}[2]{\begin{tcolorbox}[enhanced,colback=#1!11,colframe=#1!45!black,boxrule=0.42pt,arc=2mm,left=1mm,right=1mm,top=0.45mm,bottom=0.45mm,before skip=0mm,after skip=0.85mm,fontupper=\bfseries\scriptsize,colupper=#1!55!black]\centering #2\end{tcolorbox}}
\newtcolorbox{panel}[1]{colback=fondpage,colframe=fondpage,boxrule=0pt,arc=0mm,left=0mm,right=0mm,top=0mm,bottom=0mm,height=168mm,valign=top,before skip=0mm,after skip=0mm}
\newcommand{\titrepage}{\begin{tcolorbox}[colback=white,colframe=bleu!70!black,boxrule=0.65pt,arc=3mm,left=2.5mm,right=2.5mm,top=1.15mm,bottom=1.15mm,before skip=0mm,after skip=1.4mm,drop fuzzy shadow=black!10]
{\Large\bfseries Parcours de révision -- Suites}\end{tcolorbox}}
\tikzset{tabouter/.style={draw=tableline, line width=0.85pt},tabline/.style={draw=tableline, line width=0.55pt},forbidden/.style={draw=tableline, line width=0.95pt},vararrow/.style={-{Latex[length=2.7mm,width=2.0mm]}, line width=0.85pt, draw=accent},tabtext/.style={font=\small},tabmath/.style={font=\small},labelcell/.style={font=\small\bfseries, text=deepblue},pointvalue/.style={font=\small\bfseries},endvalue/.style={font=\small}}
\lstset{language=Python,basicstyle=\ttfamily\small,numbers=left,numberstyle=\tiny,frame=single,showstringspaces=false,columns=fullflexible,keepspaces=true}

\begin{document}
\begin{center}
{\LARGE\bfseries Parcours de révision -- Suites}\\[1mm]

\end{center}
\vspace{2mm}
\tableofcontents
\clearpage

\section*{Partie 1 -- Corrections des questions flashs}
\addcontentsline{toc}{section}{Partie 1 -- Corrections des questions flashs}
\setcounter{reponse}{0}
\begin{blocreponses}{Réponses -- A. Vocabulaire et formules de base}
  \item C'est une fonction définie sur \(\mathbb{N}\) ou sur une partie de \(\mathbb{N}\).
  \item On le note en général \(u_n\).
  \item Une définition explicite donne \(u_n\) directement en fonction de \(n\). Une définition par récurrence donne un terme initial et une relation reliant deux termes.
  \item Lorsqu'il existe un réel \(r\) tel que, pour tout \(n\), \(u_{n+1}=u_n+r\).
  \item On l'appelle la raison.
  \item \(u_n=u_0+nr\).
  \item Lorsqu'il existe un réel \(q\) tel que, pour tout \(n\), \(u_{n+1}=q u_n\).
  \item On l'appelle la raison.
  \item \(u_n=u_0q^n\).
  \item On compare en général \(u_{n+1}\) et \(u_n\), souvent en étudiant le signe de \(u_{n+1}-u_n\).
  \item On montre que \(u_{n+1}-u_n\geqslant 0\) pour tout \(n\).
  \item On montre que \(u_{n+1}-u_n\leqslant 0\) pour tout \(n\).
  \item Quand tous les termes de la suite sont strictement positifs.
  \item On pose souvent une suite auxiliaire \(v_n=u_n-\ell\), où \(\ell\) est un point fixe.
  \item C'est un réel \(\ell\) tel que \(\ell=a\ell+b\).
  \item Une boucle \texttt{while}.
\end{blocreponses}

\begin{blocreponses}{Réponses -- B. Limites et convergence}
  \item Cela signifie que les termes de la suite se rapprochent de \(\ell\) quand \(n\) devient très grand.
  \item Oui.
  \item Non, elle diverge vers \(+\infty\).
  \item \(0\).
  \item \(+\infty\).
  \item \(1\).
  \item Si \(u_n\leqslant v_n\leqslant w_n\) et si \((u_n)\) et \((w_n)\) ont la même limite \(\ell\), alors \((v_n)\) a aussi pour limite \(\ell\).
  \item Quand elle est majorée.
  \item Quand elle est minorée.
\end{blocreponses}

\begin{blocreponses}{Réponses -- C. Calculs flash}
  \item \(u_5=2\times 5+1=11\).
  \item \(u_1=7+2=9\).
  \item On pense à poser \(v_n=u_n-30\).
  \item Le point fixe vaut \(8\), car \(\ell=0{,}5\ell+4\).
\end{blocreponses}

\begin{blocreponses}{Réponses -- D. Suites bornées, monotones et convergentes}
  \item Une suite dont tous les termes sont inférieurs ou égaux à un même réel.
  \item Une suite dont tous les termes sont supérieurs ou égaux à un même réel.
  \item Une suite à la fois majorée et minorée.
  \item Une suite qui est toujours croissante ou toujours décroissante.
  \item Une suite dont tous les termes sont égaux.
  \item Une suite qui admet une limite finie.
\end{blocreponses}

\begin{blocreponses}{Réponses -- E. Rédaction de la récurrence}
  \item Il faut nommer clairement la propriété et préciser l'ensemble des entiers concernés : « Pour tout entier naturel \(n\), on note \(P(n)\) la propriété ... ».
  \item On vérifie la propriété au premier rang concerné, puis on le rédige explicitement : « Pour \(n=0\), ... donc \(P(0)\) est vraie. »
  \item On fixe un entier \(n\), on suppose \(P(n)\) vraie, puis on montre que \(P(n+1)\) est vraie.
  \item On rappelle que l'initialisation et l'hérédité ont été établies, puis on conclut : « Par récurrence, \(P(n)\) est vraie pour tout entier naturel \(n\) concerné. »
\end{blocreponses}

\begin{blocreponses}{Réponses -- F. Un calcul direct pour s'entraîner}
  \item \(u_{12}=u_0+12r=3+12\times 4=51\).
\end{blocreponses}


\clearpage
\section*{Partie 2 -- Corrigés des sujets de bac}
\addcontentsline{toc}{section}{Partie 2 -- Corrigés des sujets de bac}
\titreSujet{Sujet 1 -- Suite auxiliaire et algorithme}{Amérique du Nord, 21 mai 2025, jour 1, exercice 2}

\begin{correctionbox}{1. Calcul de $u_1$}
\[
u_1=\dfrac{2u_0+1}{u_0+2}=\dfrac{2\times2+1}{2+2}=\dfrac54.
\]
\end{correctionbox}

\begin{correctionbox}{2. Suite auxiliaire $(a_n)$}
\begin{enumerate}
\item On a
\[
a_0=\dfrac{u_0}{u_0-1}=2,
\qquad
a_1=\dfrac{\frac54}{\frac54-1}=5.
\]

\item Pour tout entier naturel $n$,
\[
a_{n+1}=\dfrac{u_{n+1}}{u_{n+1}-1}
=\dfrac{\frac{2u_n+1}{u_n+2}}{\frac{2u_n+1}{u_n+2}-1}
=\dfrac{2u_n+1}{u_n-1}.
\]
Or
\[
3a_n-1=3\dfrac{u_n}{u_n-1}-1=\dfrac{2u_n+1}{u_n-1}.
\]
Donc
\[
a_{n+1}=3a_n-1.
\]

\item Montrons par récurrence que, pour tout $n\geqslant1$, $a_n\geqslant3n-1$.
Pour $n=1$, $a_1=5$ et $3\times1-1=2$, donc la propriété est vraie.
Supposons $a_n\geqslant3n-1$ pour un entier $n\geqslant1$. Alors
\[
a_{n+1}=3a_n-1\geqslant3(3n-1)-1=9n-4.
\]
Or $9n-4\geqslant3n+2=3(n+1)-1$ pour $n\geqslant1$. La propriété est donc héréditaire.
Ainsi, pour tout $n\geqslant1$,
\[
a_n\geqslant3n-1.
\]

\item Comme $3n-1\to+\infty$ et $a_n\geqslant3n-1$, on obtient par comparaison
\[
\lim_{n\to+\infty}a_n=+\infty.
\]
\end{enumerate}
\end{correctionbox}

\begin{correctionbox}{3. Limite de $(u_n)$}
\begin{enumerate}
\item À partir de $a_n=\dfrac{u_n}{u_n-1}$, on obtient
\[
a_n(u_n-1)=u_n,
\]
donc
\[
u_n(a_n-1)=a_n.
\]
Ainsi
\[
u_n=\dfrac{a_n}{a_n-1}.
\]

\item On écrit
\[
u_n=\dfrac{a_n}{a_n-1}=\dfrac{1}{1-\frac1{a_n}}.
\]
Comme $a_n\to+\infty$, on a $\frac1{a_n}\to0$, donc
\[
\lim_{n\to+\infty}u_n=1.
\]
\end{enumerate}
\end{correctionbox}

\begin{correctionbox}{4. Algorithme}
\begin{enumerate}
\item Comme $(u_n)$ est décroissante et converge vers $1$, la fonction \texttt{algo(p)} renvoie le premier rang $n$ pour lequel
\[
u_n-1\leqslant p,
\]
et la valeur correspondante de $u_n$.

\item Pour $p=0{,}001$, la valeur renvoyée pour $n$ est
\[
\boxed{6}.
\]
\end{enumerate}
\end{correctionbox}

\newpage
\titreSujet{Sujet 2 -- Deux modèles d'évolution d'une population}{Centres étrangers, 13 juin 2025, sujet 2, exercice 1}

\begin{correctionbox}{Partie A}
\begin{enumerate}
\item Au 1er janvier 2026, on a
\[
u_1=6\times0{,}93=5{,}58.
\]
La population est donc de $5\,580$ individus.

\item La suite $(u_n)$ est géométrique de premier terme $6$ et de raison $0{,}93$, donc
\[
u_n=6\times0{,}93^n.
\]

\item Comme $0<0{,}93<1$,
\[
\lim_{n\to+\infty}0{,}93^n=0,
\]
donc
\[
\lim_{n\to+\infty}u_n=0.
\]
Selon ce modèle, la population du milieu $A$ tend vers $0$ millier d'individus.
\end{enumerate}
\end{correctionbox}

\begin{correctionbox}{Partie B}
\begin{enumerate}
\item On a
\[
v_1=-0{,}05\times6^2+1{,}1\times6=-1{,}8+6{,}6=4{,}8.
\]
La population est donc de $4\,800$ individus au 1er janvier 2026.

\item Pour $x\geqslant0$,
\[
f'(x)=-0{,}1x+1{,}1.
\]
Ainsi $f'(x)\geqslant0$ sur $[0;11]$. La fonction $f$ est donc croissante sur $[0;11]$.

\item Montrons par récurrence que $2\leqslant v_{n+1}\leqslant v_n\leqslant6$.
Pour $n=0$, $v_0=6$ et $v_1=4{,}8$, donc $2\leqslant v_1\leqslant v_0\leqslant6$.
Supposons $2\leqslant v_{n+1}\leqslant v_n\leqslant6$. Comme $f$ est croissante sur $[0;11]$,
\[
f(2)\leqslant f(v_{n+1})\leqslant f(v_n).
\]
Or $f(2)=2$ et $f(v_{n+1})=v_{n+2}$, $f(v_n)=v_{n+1}$. Donc
\[
2\leqslant v_{n+2}\leqslant v_{n+1}\leqslant6.
\]
La propriété est héréditaire.

\item La suite $(v_n)$ est décroissante et minorée par $2$, donc elle converge vers une limite $\ell$.

\item La fonction $f$ est continue, donc en passant à la limite dans $v_{n+1}=f(v_n)$, on obtient $\ell=f(\ell)$.
Ainsi
\[
\ell=-0{,}05\ell^2+1{,}1\ell.
\]
Donc
\[
-0{,}05\ell^2+0{,}1\ell=0,
\]
soit
\[
\ell(-0{,}05\ell+0{,}1)=0.
\]
Les solutions sont $0$ et $2$. Comme $v_n\geqslant2$, on a
\[
\ell=2.
\]
Selon ce modèle, la population du milieu $B$ tend vers $2\,000$ individus.
\end{enumerate}
\end{correctionbox}

\begin{correctionbox}{Partie C}
\begin{enumerate}
\item On résout
\[
6\times0{,}93^n<3.
\]
Cela équivaut à
\[
0{,}93^n<\dfrac12.
\]
Avec les logarithmes :
\[
n\ln(0{,}93)<\ln\left(\dfrac12\right).
\]
Comme $\ln(0{,}93)<0$, on obtient
\[
n>\dfrac{\ln(1/2)}{\ln(0{,}93)}\approx9{,}55.
\]
Le plus petit entier convenable est $n=10$, donc l'année est
\[
2025+10=\boxed{2035}.
\]

\item À la calculatrice, on obtient $v_6\approx2{,}965<3$ alors que $v_5\approx3{,}145>3$. L'année cherchée est donc
\[
2025+6=\boxed{2031}.
\]

\item On a $u_n\to0$ et $v_n\to2$. Donc, à partir d'un certain rang, $u_n$ est proche de $0$ et $v_n$ est proche de $2$. La population du milieu $B$ dépassera donc celle du milieu $A$ à partir d'une certaine année.

\item Un programme possible est :
\begin{Verbatim}[frame=single, fontsize=\small]
n = 0
u = 6
v = 6
while v <= u:
    u = 0.93*u
    v = -0.05*v**2 + 1.1*v
    n = n + 1
print(2025 + n)
\end{Verbatim}
À l'exécution, l'année affichée est
\[
\boxed{2038}.
\]
\end{enumerate}
\end{correctionbox}

\newpage
\titreSujet{Sujet 3 -- Convergence de deux suites vers une même limite}{Métropole, 9 septembre 2025, sujet 1, exercice 3}

\begin{correctionbox}{Partie A}
\begin{enumerate}
\item La fonction $f$ est définie sur $[2;+\infty[$ par $f(x)=\sqrt{3x-2}$. Elle est dérivable sur $]2;+\infty[$ et
\[
f'(x)=\dfrac{3}{2\sqrt{3x-2}}>0.
\]
Elle est donc croissante. De plus $f(2)=2$ et $\lim_{x\to+\infty}f(x)=+\infty$.

\begin{center}
\begin{tikzpicture}[x=1cm,y=1cm]
\def\L{2.2}\def\W{8.6}\def\H{2.55}\def\Yx{1.75}
\fill[softgray] (0,\Yx) rectangle (\W,\H);
\fill[softblue] (0,0) rectangle (\L,\Yx);
\draw[tabouter] (0,0) rectangle (\W,\H);
\draw[tabline] (0,\Yx) -- (\W,\Yx);
\draw[tabline] (\L,0) -- (\L,\H);
\node[labelcell] at ({\L/2},{(\Yx+\H)/2}) {$x$};
\node[labelcell] at ({\L/2},{\Yx/2}) {$f(x)$};
\node[pointvalue] at ({\L+0.35},{(\Yx+\H)/2}) {$2$};
\node[tabmath] at ({\W-0.55},{(\Yx+\H)/2}) {$+\infty$};
\node[pointvalue] at ({\L+0.45},0.35) {$2$};
\node[endvalue] at ({\W-0.75},1.30) {$+\infty$};
\draw[vararrow] ({\L+0.9},0.5) -- ({\W-1.2},1.25);
\end{tikzpicture}
\end{center}

\item On utilise le fait que, pour $x\geqslant2$,
\[
f(x)\leqslant x
\]
car $\sqrt{3x-2}\leqslant x \Longleftrightarrow 3x-2\leqslant x^2 \Longleftrightarrow (x-1)(x-2)\geqslant0$.
On montre alors par récurrence que $2\leqslant u_{n+1}\leqslant u_n\leqslant6$. La propriété est vraie au rang $0$ car $u_0=6$ et $u_1=4$.
Si elle est vraie au rang $n$, alors $u_{n+1}\geqslant2$, donc
\[
2\leqslant f(u_{n+1})\leqslant u_{n+1}.
\]
Or $f(u_{n+1})=u_{n+2}$, d'où
\[
2\leqslant u_{n+2}\leqslant u_{n+1}\leqslant6.
\]
La suite $(u_n)$ est donc décroissante et minorée par $2$ ; elle converge.

\item Si $\ell$ est la limite de $(u_n)$, on admet que $f(\ell)=\ell$. Ainsi
\[
\sqrt{3\ell-2}=\ell.
\]
Comme $\ell\geqslant2$, on peut élever au carré :
\[
3\ell-2=\ell^2,
\]
donc
\[
(\ell-1)(\ell-2)=0.
\]
Puisque $\ell\geqslant2$, on obtient
\[
\ell=2.
\]

\item Comme $(u_n)$ converge vers $2$ en restant au-dessus de $2$, l'appel \texttt{rang(2.000001)} renvoie une valeur. Plus généralement, \texttt{rang(a)} renvoie un résultat exactement lorsque
\[
a>2.
\]
Si $a\leqslant2$, la condition $u\geqslant a$ reste toujours vraie.
\end{enumerate}
\end{correctionbox}

\begin{correctionbox}{Partie B}
\begin{enumerate}
\item On calcule
\[
v_1=3-\dfrac2{6}=\dfrac83.
\]

\item On a
\[
w_{n+1}=\dfrac{v_{n+1}-1}{v_{n+1}-2}
=\dfrac{2-\frac2{v_n}}{1-\frac2{v_n}}
=\dfrac{2v_n-2}{v_n-2}
=2\dfrac{v_n-1}{v_n-2}=2w_n.
\]
La suite $(w_n)$ est donc géométrique de raison $2$ et
\[
w_0=\dfrac{6-1}{6-2}=\dfrac54=1{,}25.
\]
Donc
\[
w_n=1{,}25\times2^n.
\]
Comme $w_n-1=\dfrac1{v_n-2}$, on obtient
\[
v_n=2+\dfrac1{w_n-1}=2+\dfrac{1}{1{,}25\times2^n-1}.
\]
Ainsi
\[
\lim_{n\to+\infty}v_n=2.
\]

\item On résout
\[
2+\dfrac{1}{1{,}25\times2^n-1}<2{,}01.
\]
Cela équivaut à
\[
\dfrac{1}{1{,}25\times2^n-1}<0{,}01,
\]
donc
\[
1{,}25\times2^n-1>100.
\]
Ainsi
\[
2^n>80{,}8.
\]
Or $\log_2(80{,}8)\approx6{,}34$. Le plus petit entier convenable est donc
\[
\boxed{n=7}.
\]
\end{enumerate}
\end{correctionbox}

\begin{correctionbox}{Partie C}
D'après la partie B, $v_n<2{,}01$ à partir du rang $7$ et $v_n>2$ pour tout $n$, donc $v_n\in]1{,}99;2{,}01[$ pour tout $n\geqslant7$.
Pour la suite $(u_n)$, la calculatrice ou la fonction \texttt{rang} donne
\[
u_{17}\approx2{,}0087<2{,}01
\]
et $u_n\geqslant2$ pour tout $n$. Ainsi $u_n\in]1{,}99;2{,}01[$ pour tout $n\geqslant17$.
Le plus petit entier convenable est donc
\[
\boxed{N=17}.
\]
\end{correctionbox}


\end{document}
